\section{Protocol Pseudocode}
\label{app:algorithms}
These six procedures specify the normative protocol, not the current CLI.
They consume authenticated predecessors; they do not accept a caller's
assertion that an omitted stage happened. $\Sigma$ denotes authoritative
registry/trust/budget state; $J$ denotes the lifecycle record store.
The event types and permitted edges are those of Table~\ref{tab:lifecycle}.
For a government instance, the bound profile $P_G$ and stage-due predicates
$G^k$ are those of Section~\ref{sec:gov-profile}; they strengthen existing
guards without creating new phases or JSON schemas.
Every successful progression destination checks the current $M_G^k(t)$ at authoritative
transition time, not the due set captured during an earlier successful call.
Denial and other no-service transitions keep their original authority checks;
they do not require passing institutional obligations.
A \emph{refusal} preserves a failure record, emits no new authorization, and
returns BLOCK as the recommendation for missing/invalid evidence.
Its lifecycle state is ABORTED for integrity/conflict/stale-head failure,
REDESIGN\_REQUIRED for an incomplete/remediable proposal, or no service for a
serving failure. Legacy internal gate vocabulary is not a facade disposition.

\subsection{Algorithm A1: VerifyEnvelope$(M,I,q,J,\Sigma)$}
\begin{enumerate}
\item Capture the exact envelope and referenced source bytes. Reject unknown
version/type, malformed or noncanonical encoding and invalid field ranges.
Hash and parse the same captured bytes, not two reads of mutable paths.
\item Check the domain-separated signed envelope, live issuer key, role/type
permission, instance digest, issue/expiry times and required trust profile.
A caller-supplied public key is not by itself an approved authority.
\item If message ID was seen with different digest, abort. If identical,
return the recorded historical receipt with \emph{no new state transition}.
Historical receipt reuse does not certify current serving eligibility.
\item Require phase $q$ and the exact allowed edge; load every required
predecessor by committed identifier. Check its digest, signature, instance,
continuing validity and semantic linkage. Reject absence or substitution.
\item Replay payload-specific semantics using captured inputs. Return the
validated payload and captured digests, never a bare signature-success flag.
\end{enumerate}

\subsection{Algorithm A2: Freeze$(R,P,N,\Sigma)$}
\begin{enumerate}
\item Verify owner Registration from DRAFT with A1. Require current head,
policy/trust/population snapshots, legitimate purpose, accountable owner,
recipients, artifacts, complete interface, candidates and governance context.
For a government instance, bind the competent authority's approved $P_G$,
delegations, intended/prohibited uses, affected population and mandatory
institutional roster; require $G^{\rm register}$.
Persist REGISTERED bound to immutable $I$.
\item Register each requested claim's approved O1--O7 witness route and failure
action. A scoped no-go annotation binds its observation class and premises;
it never waives evidence or supplies a safety certificate.
Check finite scenario partition, nonempty in-scope threat sets, channel
and joint-observation witnesses, exclusions and utility/search requirements.
Freeze the complete evidence roster and allowed evidence directions.
For $P_G$, freeze stage-due owners, bounded checkers, review/expiry conditions,
oversight/redress arrangements and approved disclosure exemptions. Later-due
judgments are scheduled obligations, not assumed completed reviews.
\item Obtain EvidencePlan approval by the assessor and required approvers.
Require bound workers, training/source provenance obligations, positive
controls, exclusions, sampling/selection/stopping rules and failure retention.
\item Before covered observations, reserve the family's statistical allocation
in the authoritative lifetime ledger. Include related failed/retried attempts;
require a justified simultaneous procedure, not merely a numerical budget.
\item Verify the approved plan envelope, including $G^{\rm plan}$ where
applicable, and persist PLAN\_FROZEN. Any changed
instance field requires a new identifier and applicable stages again.
\end{enumerate}

\subsection{Algorithm A3: CollectAndAssess$(I,J)$}
\begin{enumerate}
\item Require PLAN\_FROZEN. Run the registered workers against bound sources;
retain successes, failures and exclusions. Verify each EvidenceBundle with
A1, including observed provenance, controls and required training/completion
records. Missing evidence refuses; worker signatures do not prove truth.
Collect profile-required institutional evidence assigned to these workers;
keep its review status separate from the numerical privacy endpoints.
\item Close the exact bundle inventory. For government instances require
$G^{\rm evidence}$ before persisting EVIDENCE\_FROZEN.
For each candidate and every in-scope threat, replay evidence directions,
game/population identity, arithmetic and admissibility.
\item Call Admit under O1--O5/O7 for every cell; require every prescribed
source, control, scope and trust witness. Refuse missing/failed required
routes, recording the precise reason and remedy or scoped obstruction.
When no nontrivial upper endpoint is available, only an explicitly allowed
worst-case ceiling with complete required evidence may return Valid rather
than Refuse. Run each approved bounded route checker; timeout, an unsupported
route or failed check refuses. Derive $L,U$ only from compatible jointly
covered evidence. Verify direct
release binding or a safe-direction transfer certificate with its error
penalty. Check the full joint portfolio at the expected head. Preserve
unsupported or contradictory inputs as failures, not zero observations.
\item Set internal cell status to BLOCK if a valid $L>\tau$, CLEAR only when
all cell obligations pass and $U\le\tau$, otherwise INCONCLUSIVE for genuine
unresolved valid evidence. Invalid/missing obligations yield refusal at the
facade even if a lower-level gate uses another vocabulary.
\item Sign the complete AssessmentReport including failures and assumptions;
for government instances require $G^{\rm assess}$ before persisting ASSESSED.
A report is not authorization and cannot certify live service.
\end{enumerate}

\subsection{Algorithm A4: RecommendAndSelect$(I,A,Z)$}
\begin{enumerate}
\item Require ASSESSED; apply O6 and revalidate assessment $A$, original request/scope and
optional risk acceptance $Z$. For each candidate form $B$ and $X$ exactly as
in Section~\ref{sec:theory-protocol}; include missing in-scope evidence and
blocking floors regardless of mandatory/optional labels.
Validate any resolving plan against the instance, estimand, procedure,
obtainable data, unspent allocation and decision-changing target. An advisory
action string or a budget already consumed by inspected evidence is not a
valid plan. Portfolio conditions below cover $H$ plus the candidate, not only
the active service roster.
\item \textbf{If} $B\ne\varnothing$, set $d=$BLOCK.
\textbf{Else if} $X=\varnothing$, set $d=$RELEASE.
\textbf{Else if} $Z$ is valid and binds every crossing, set $d=$RELEASE-WITH-RISK.
\textbf{Else if} every crossing has a valid resolving plan, set $d=$INCONCLUSIVE.
\textbf{Else} set $d=$BLOCK. Emit the bound per-threat record.
\item Build $F$ only from dispositions RELEASE or RELEASE-WITH-RISK that also
pass every mandatory CLEAR, utility, control, transfer, portfolio and policy
condition. RWR requires every crossing to be optional, policy-permitted and
exactly accepted. BLOCK and INCONCLUSIVE are ineligible. Acceptance cannot
replace a mandatory CLEAR.
For government instances also require $G^{\rm select}$; FAIL or UNKNOWN
excludes the candidate without converting the institutional judgment into
a privacy score. Independent authorization review remains due in A5.
\item If $F$ is empty, return REDESIGN\_REQUIRED, unless an approved complete
candidate-search certificate proves infeasibility; only then return REJECTED.
Neither refusal implies that an unrestricted alternative model is impossible.
\item Otherwise choose the candidate by the frozen objective/tie-break,
retaining every candidate's result. Sign OptimizationReport bound to
the complete portfolio and persist OPTIMIZED. A recommendation or selection
alone cannot start A6.
\end{enumerate}

\subsection{Algorithm A5: ReviewAndCommit$(I,O,J,\Sigma)$}
\begin{enumerate}
\item Require OPTIMIZED. Revalidate assessment, recommendation, selection $O$,
policy, populations, evidence and all mandatory G0--G11 obligations. Replay
A4's disposition eligibility, including the optional-crossing acceptance rule;
a mandatory-only test cannot admit an INCONCLUSIVE candidate.
Require the authorizer's reasoned decision, independent challenge,
conflict/objection dispositions, conditions, expiry and incident/retirement
ownership. An RWR record cannot discharge an unresolved mandatory gate.
For $P_G$, obtain the competent independent institutional reviews now due;
require $G^{\rm review}$ with current reasons, conditions and delegation status.
A privacy RELEASE cannot replace a mandate, rights or redress obligation.
\item Verify AuthorizationCommitRequest; recheck $G^{\rm review}$ where
applicable before recording it and entering COMMIT\_PENDING.
Within one registry transaction require expected head, unused nonce and
unchanged selection. Recheck live evidence, trust, revocations, full current
portfolio over the cumulative disclosure history $H$, including retired and
revoked releases, and availability of every cumulative budget delta. Check
history continuity across rebasing or changes of purpose; a new instance does
not establish a new adversary. Reserve the candidate's possible exposure in
the history if the authorization transaction itself discloses it.
For government instances recheck GovCommitOK against the same policy,
delegation and review-status snapshot and bind that snapshot to the transition.
\item Compute the complete transition digest and next head; prepare the
authorization. Atomically persist transition, nonce, budget deltas,
authorization and head. On failure abort without exposing a usable receipt.
\item After successful commit publish AuthorizationReceipt and independently
verifiable inclusion/consistency evidence; enter AUTHORIZED.
A stale retry rebases into a new instance and replays affected obligations.
Already inspected statistical evidence is not refunded.
\end{enumerate}

\subsection{Algorithm A6: ActivateServeMonitor$(I,\Sigma)$}
\begin{enumerate}
\item Require AUTHORIZED; enforce O7 through authentic current registry inclusion/status.
Remeasure served artifact, complete interface and controls; compare exact
bindings, verify current evidence/trust and issue a policy-bounded lease.
Require $G^{\rm serve}$ for government instances, including the operational
oversight, redress, intended-use and renewal conditions currently due.
Only then record ActivationReceipt and enter ACTIVE.
\item \textbf{For every new access:} capture the immutable measured
artifact/interface snapshot. At one service linearization point serialized
with revocation, validate ACTIVE, current authorization epoch/inclusion,
snapshot identity, live lease/evidence, allowed query and no drift/incident.
For government instances also require current $G^{\rm serve}$ against the
same authoritative status/epoch; unmet, stale or unknown institutional
conditions prevent a grant even when privacy remains CLEAR.
Jointly reserve required budgets, append any newly exposed release to $H$,
and grant access to that snapshot. Uncertain completion cannot remove an
exposure reservation. A stale
epoch or failed budget/status check aborts the grant.
\item \textbf{On unknown status or failed check:} deny access before or
atomically with suspension recording. Expiry yields EXPIRED; authorized
revocation yields REVOKED. Preserve incidents and require a fresh applicable
assessment path for changed or suspended service.
Profile renewal deadlines, withdrawal of authority and retirement conditions
use these same no-service transitions; a new purpose or mandate cannot reuse
the former instance's authorization.
\item Independently monitor registry consistency, gateway conformance,
interface changes and incidents. Never use a cached historical signature
to bypass the per-access guard. Revocation prevents grants linearized after
it; completion of an earlier grant is not ruled out. Disclosed information
is not recalled. Preserve $H$ after expiry, revocation and retirement. These
terminal service states do not terminate applicable records, redress or
correction obligations; keep an accountable successor owner and an authorized
retention/access policy. Stopping a model-backed public service invokes the
approved continuity or human-service fallback, not automatic abandonment of
the underlying public function.
\end{enumerate}

\section{Proofs of the Protocol Guarantees}
\label{app:proofs}
All proofs in this appendix are mathematical arguments about the specified
objects. They are not new Lean kernel checks or verified Python refinements.

\subsection{Proof of Proposition~\ref{prop:decision}}
\begin{proof}
The ordered branches are mutually exclusive after precedence is applied and
exhaust all inputs, including the final refusal branch. RELEASE excludes $B$
and $X$, hence every valid interval has $U_t\le\tau_t$.
On joint coverage, $\theta_t\le U_t\le\tau_t$ for every in-scope threat.
A blocking floor gives $\theta_t\ge L_t>\tau_t$. A crossing contains values
on both sides of tolerance and establishes neither assertion by endpoints
alone. Missing/invalid evidence enters $B$ irrespective of endpoint displays
or acceptance; thus an administrative BLOCK need not demonstrate high risk.
Acceptance has no influence on the earlier blocking branch.
\end{proof}

\subsection{Proof of Theorem~\ref{thm:statistics}}
\begin{proof}
On $\bigcap_j E_j$, all upper bounds selected at any covered opportunity are
valid, regardless of selection's dependence on covered observations.
Sound deterministic ceilings remain valid as well.
Proposition~\ref{prop:decision} excludes any false ordinary RELEASE there.
Thus the false-release event is a subset of $\bigcup_jE_j^c$ and
\[
\begin{aligned}
 \Pr\Big(\bigcup_jE_j^c\Big)
 &\le \sum_j\mathbb E[\Pr(E_j^c\mid\mathcal F_{j-1})]\\
 &\le \mathbb E\Big[\sum_j\alpha_j\Big]\le\alpha .
\end{aligned}
\]
For countably many families take increasing finite unions and use monotone
convergence. Independence between families is not used.
A failed/aborted family capable of influencing later choices belongs in this
union. Reusing exactly covered evidence adds no new failure event, whereas
new uncovered looks require fresh valid coverage. The argument assumes
scientific adequacy and correct reduction; it does not assign their failures
probability zero or include RWR as below-threshold clearance.
\end{proof}

\subsection{Proof of Proposition~\ref{prop:scope}}
\begin{proof}
Set equality establishes coverage; pairwise disjointness establishes uniqueness.
For predicates $P_1,\ldots,P_m$ that are total on a finite universe, each element
has exactly one Boolean profile, hence exactly one atom
$A_b=\{x:(P_i(x))_i=b\}$. Omitting the all-false atom fails equality whenever
it is nonempty. Neither proof establishes that these predicates or this
universe describe all real attackers. The hidden-channel worlds in
Appendix~\ref{app:gaps}, together with Theorem~\ref{thm:observations}, disprove
that stronger implication for observation-limited verifiers.
\end{proof}

\subsection{Proof of Theorem~\ref{thm:transfer}}
\begin{proof}
For a released decoder $D$, compose $Q$ and $D$. When $K_r=K_aQ$, the
composition gives exactly the same conditional action law under $K_a$.
Maximization over assessed decoders yields the claimed order.
For approximate transfer let
$f_s(y)=\sum_aD(a\mid y)g(s,a)\in[0,1]$.
For each state, its expected value under $K_r(\cdot\mid s)$ differs from
that under $(K_aQ)(\cdot\mid s)$ by at most row TV, hence by $\eta$.
Average with $\pi$, maximize over $D$, then intersect with the trivial
ceiling one. The composed decoder must be allowed by the adversary class;
without that closure or with a different loss metric the proof does not apply.
All observations used by $D$, including prior releases and side channels,
must be in the channel being compared.
\end{proof}

\subsection{Proof of Theorem~\ref{thm:lifecycle}}
\begin{proof}
Induct on successful transitions from DRAFT. The base state serves nothing.
A1 admits only a role-authorized edge with all exact, live predecessors.
A2 binds O1--O7 witness routes before observation; A3 refuses missing routes,
invalid direction, uncovered selection, absent joint evidence and unsupported
metric transfer. A4 propagates refusals and preserves blocking precedence;
A5 cannot use acceptance to waive mandatory CLEAR. A6 supplies the separate
live-binding guard. A no-go annotation changes none of these acceptance
predicates. Thus dropping a witness does not become a successful transition.
Each successive pre-commit procedure adds its required record without changing
$I$. Invalid events add no successful progression. Identical-message retry
adds no edge; conflicting content is rejected. Therefore any COMMIT\_PENDING
state has the ordered mandatory predecessor chain and bound selection.

A5 is the only authorizing transition. Linearizability gives one successful
comparison against the expected head, and atomic persistence makes the
authorization inseparable from its portfolio/budget update. Two races against
the same head cannot both commit that transition: the first changes the head,
so the second aborts. An uncommitted signature is not sufficient for A6,
which checks authoritative inclusion and current status.

A6 is the only serving path. Its service linearization point checks committed
live authorization, exact immutable measured bindings, allowed budget and
absence of suspension, expiry or revocation while reserving budget and granting
access atomically. Serialization excludes a revocation between the checked
status and a later grant. If revocation linearizes first, the status/epoch
check fails; if the grant linearizes first it is an already granted request,
not post-revocation admission. A failed or unknown check denies access. Thus
every served request has the stated integrity properties. This is a safety
argument under faithful transition evaluation and complete mediation, not a
proof that a host cannot serve through an unmodeled path.

If service also requires ordinary RELEASE with matching current context, every
served threshold violation is a false RELEASE in
Theorem~\ref{thm:statistics}'s event. The probability bound follows by inclusion.
This bound assumes the stated scientific and infrastructure premises in the
same probability model as the evidence guarantees. No unconditional numerical
bound is claimed when those premises fail. Accepted optional residual risk is not
silently included in this ordinary-release guarantee.

For conditional progress, the principal path has eight successful edges to
ACTIVE. Induct on the remaining edges: eventual successful responses and
terminating guards advance the next edge; no conflicting external head change
before the instance's own commit and live deadlines
prevent abort/expiry. After finitely many such advances A6 activates.
Withheld messages, repeated conflicting head changes or expiring evidence
invalidate those hypotheses. A single valid Lean trace establishes
nonvacuity, not this progress statement or distributed availability.
\end{proof}

\paragraph{Government-profile restriction}
Ignoring additional institutional annotations, every successful transition
under $P_G$ is an existing MRAP transition: the predicates $G^k$ only remove
permitted successes. At each transition they use the then-current due set
$M_G^k(t)$ and authoritative status; advancing time or an arrived review
deadline cannot reuse an earlier PASS without the required current evidence.
Induction therefore preserves the conditional lifecycle
safety invariant. The ordinary-release probability bound also remains valid
when all government-influenced choices remain inside its covered family;
it does not bound legal, fairness or institutional failure probabilities.
Conditional progress additionally requires the stronger institutional guards
to be satisfiable, with competent responses and live reviews. This argument
proves no substantive institutional judgment from its attestation.

\subsection{Proof of Theorem~\ref{thm:finite-ceiling}}
\begin{proof}
Write $\mu_a$ and $\widehat\mu_a$ for true and empirical gain of decoder $a$.
For each fixed decoder, Hoeffding gives
$\Pr(|\widehat\mu_a-\mu_a|>e_n)\le2\exp(-2ne_n^2)=\alpha/M$.
Union over all $M$ decoders gives simultaneous accuracy at least $1-\alpha$.
Shared samples across decoders do not require independence between these
error events. On that event,
\[
 \max_a\mu_a-e_n\le h\le\max_a\mu_a+e_n.
\]
Randomized decoders cannot improve a linear objective beyond a deterministic
one, so completeness gives $\max_a\mu_a=\theta$.
Intersect with $[b,1]$ to obtain the reported interval.
An inconsistent interval off the coverage event must be rejected, not sorted
or silently repaired.

On the same event $U\le\theta+2e_n$ and $L\ge\theta-2e_n$.
If $\theta\le\tau-\delta$ and $2e_n<\delta$, then $U<\tau$; if
$\theta\ge\tau+\delta$, then $L>\tau$.
Solving $2e_n<\delta$ gives the sufficient sample size.
There are $|A|$ independent action choices for each of $|Y|$ observations,
giving $M=|A|^{|Y|}$. Nevertheless the empirical maximization separates:
\[
 h=\frac1n\sum_y\max_a\sum_{i:Y_i=y}g(S_i,a).
\]
Grouping labeled observations and evaluating all actions costs at most
$O(n|A|)$ gain evaluations, or $O(n+|S||Y||A|)$ using joint counts and a
gain table. Thus enumerating all decoders is not necessary for this complete,
unrestricted finite class. A certified finite alphabet and sufficient labeled
IID samples, potentially enormous for adaptive transcripts, are still required.
Incomplete/data-selected decoders without valid joint
coverage, non-IID sampling, or an unmodeled transcript invalidate the inference.
With a simultaneous channel set $\mathcal K$ containing the true $K$,
$\theta\le\sup_{K'\in\mathcal K}V_g(\pi,K')$ follows directly; calculating
that supremum soundly is a separate obligation.
\end{proof}

\section{Gap-Closure and Impossibility Proofs}
\label{app:gaps}
\begin{theorem}[Observation-relative nonidentifiability]
\label{thm:observations}
Let admissible worlds $W_0,W_1$ have
$\theta(W_0)\le\tau<\theta(W_1)$. Let $P_0,P_1$ be the distributions of
everything a verifier observes, including available records. Every randomized
verifier satisfies
\[
 P_1(\Release)\ge P_0(\Release)-d_{\rm TV}(P_0,P_1).
\]
Thus soundness $P_1(\Release)\le\alpha$ entails
$P_0(\Release)\le\alpha+d_{\rm TV}(P_0,P_1)$.
\end{theorem}
\begin{proof}
Given an observation $o$, write $r(o)\in[0,1]$ for the probability that the
verifier returns RELEASE, integrating its independent random coins.
The two release probabilities are $\int r\,dP_i$; their difference is at
most total variation by its bounded-function characterization. Identical
observation laws give identical release probabilities.
This restricts the chosen observation map, not every conceivable instrument.
\end{proof}

\subsection{G1: floor nonidentifiability and its repair}
\begin{proof}
Let $S$ be a uniform bit. A constant-output channel $K_0$ has guessing value
$1/2$; $K_1$ reveals $S$ and has value one. A decoder that always answers zero
has success $1/2$ in both. If the observed evidence is this success summary,
the two observation laws agree. At any $\tau\in(1/2,1)$,
Theorem~\ref{thm:observations} forbids informative sound clearance from that
summary alone. A deterministic upper bound valid for every compatible
world must be at least one. This does not apply when additional information
identifies the channel. Exact channel evaluation enumerates
$\sum_y\max_a\sum_s\pi(s)K(y\mid s)g(s,a)$; complete finite-decoder estimation
is justified by Theorem~\ref{thm:finite-ceiling}. The repair therefore adds
information/completeness assumptions, not an inversion of empirical failure.
\end{proof}

\subsection{G2: multiplicity, repair, and resolution limits}
\begin{proof}
Let every true risk be one, and let each upper bound be zero with probability
$p$ and one otherwise. Its marginal coverage is $1-p$.
With $m$ independent bounds and a tolerance below one, choosing the smallest
falsely clears with probability $1-(1-p)^m$. At $p=1/20,m=20$ this is
$1-(19/20)^{20}$. Replacing each allocation by $p=\alpha/m$ gives
$1-(1-\alpha/m)^m\le\alpha$ by the union bound. The bound remains valid
without independence if each error probability is at most $\alpha/m$.
This proves the concrete repair, not validity of arbitrary imported intervals.
An unbounded sequence at fixed $p>0$ makes the independent failure probability
tend to one. Predictable nonrefundable spending with genuine family validity
instead satisfies Theorem~\ref{thm:statistics}.
\end{proof}

\subsubsection{No uniform finite-budget resolution at vanishing margins}
\label{app:threshold}
\begin{proposition}[Fixed-budget threshold boundary]
\label{prop:threshold}
Fix $0<\tau<1$, $n\ge1$, and $0\le\alpha<1/2$. A possibly randomized rule
observes $n$ IID Bernoulli$(p)$ samples and outputs LOW, HIGH or UNRESOLVED.
Assume uniformly
$\sup_{p\le\tau}P_p(\mathrm{HIGH})\le\alpha$ and
$\sup_{p>\tau}P_p(\mathrm{LOW})\le\alpha$.
For $p_\pm=\tau\pm\varepsilon$ inside $(0,1)$,
\[
 P_{p_\pm}(\mathrm{LOW\ or\ HIGH})
 \le 2\alpha+\Delta_n
 \le 2\alpha+\min\{1,2n\varepsilon\},
\]
where $\Delta_n=d_{\rm TV}(\mathrm{Ber}(p_-)^{\otimes n},
\mathrm{Ber}(p_+)^{\otimes n})$.
\end{proposition}
\begin{proof}
Under $p_-$, LOW has probability at most its probability under $p_+$ plus
$\Delta_n$, hence at most $\alpha+\Delta_n$; HIGH has probability at most
$\alpha$. Interchange the sides for the other conclusion.
Couple each Bernoulli pair with a common uniform variate; the probability of
a mismatched pair is $2\varepsilon$.
Therefore $\Delta_n\le1-(1-2\varepsilon)^n\le2n\varepsilon$, also at most one.
For fixed $n$, each event probability is a polynomial in $p$, even with
randomized decisions. Continuity as $p\downarrow\tau$ gives
$P_\tau(\mathrm{LOW})\le\alpha$, so conclusive probability at the threshold
is at most $2\alpha$. No fixed-$n$ rule satisfying the error guarantees
has conclusive probability uniformly greater than $2\alpha$ over every
positive margin.
\end{proof}
LOW/HIGH mean scientific conclusions. Missing-evidence BLOCK and RWR count
as unresolved risk, not successful scientific resolution. This result does
not forbid increasing samples for each fixed margin, exact analytic
certificates, or useful power at a registered positive margin.

\subsection{G3: finite coverage and invisible channels}
\begin{proof}
Proposition~\ref{prop:scope} proves the finite declared check, including its
residual cell. For the unqualified claim take one world with only a declared
constant output and another with the same supplied records but an unobserved
log revealing $S$. The verifier's observation laws coincide while risks differ.
Apply Theorem~\ref{thm:observations}. Externally justified closed-world
instrumentation may distinguish the worlds; the impossibility no longer
applies when its observation premise is removed. Partition tests alone do
not verify such instrumentation or the reasons for exclusions.
\end{proof}

\subsection{G4: marginal equivalence does not imply joint equivalence}
\begin{proof}
In the safe world let $Y_1,Y_2$ be independent uniform bits independent of $S$.
In the unsafe world take independent uniform $S,R$ and
$(Y_1,Y_2)=(R,S\oplus R)$. In both worlds each conditional marginal given
$S=s$ is uniform. Marginal observation channels, not only marginal success
scores, therefore agree. The safe joint channel is independent of $S$, giving
value $1/2$; the unsafe joint pair determines $S=Y_1\oplus Y_2$, giving one.
Consequently a universal joint upper rule using only these marginals cannot
return a ceiling below one for this pair.
Full joint observation distinguishes the worlds and exact joint Bayes
evaluation recovers their values. A valid joint garbling certificate invokes
Theorem~\ref{thm:transfer}; no such conclusion follows from marginal garbling
certificates alone.
\end{proof}

\subsection{G5: TV sharpness and the posterior obstruction}
\begin{proof}
For a uniform bit, compare a channel whose binary output is uniform and independent of
$S$ with a channel that answers $S$ correctly with probability $1/2+\eta$,
$0\le\eta\le1/2$. Row TV is $\eta$ and guessing values are $1/2$ and
$1/2+\eta$, attaining the additive bound with identity $Q$.

Separately, let the assessed channel always emit $\bot$, and let the released
channel emit $S$ with probability $\varepsilon>0$, otherwise $\bot$.
Row TV is $\varepsilon$ and expected guessing is $1/2+\varepsilon/2$,
consistent with the TV theorem. Conditional on either revealing output,
posterior success is one for every positive $\varepsilon$.
Thus no uniformly nontrivial maximum-posterior ceiling follows from arbitrarily
small row TV alone. A likelihood-ratio constraint or other suitable
observation-specific premise would be additional information, not a consequence
of the expected-gain transfer theorem.
\end{proof}

\subsection{G6: acceptance noninterference}
\begin{proof}
In A4, $B$ is evaluated before acceptance. If evidence is missing, invalid,
stale, contradictory or a floor violates tolerance, all later acceptance
branches are unreachable. For a valid crossing, acceptance writes a different
disposition without changing $L,U,\tau$ or the channel. Choose any
$\theta\in(\tau,U]$ compatible with the interval: RWR can coexist with
excessive true risk. At $[0.4,0.8]$, $\tau=0.6$, both $\theta=0.5$ and
$\theta=0.7$ fit. The same interval therefore supports neither a safety proof
nor an inevitable violation. A4's strict feasibility and A5's conjunctive
gates separately exclude override of a mandatory non-CLEAR cell.
The finite-grid experiment checks a concrete reducer implementation, not
all parsing, signing, authority or state-transition behavior.
\end{proof}

\subsection{G7: authentic false accounts}
\begin{proof}
Fix an ideal authenticated record claiming interval $[1/2,11/20]$ at
$\tau=3/5$. Its signer is authorized and its bytes are unchanged.
One world implements a constant-output channel of value $1/2$; another
implements a revealing channel of value one but its signer reports the same
account. A verifier restricted to that authenticated account has identical
observations, so Theorem~\ref{thm:observations} applies.
No signature was forged and the displayed interval is internally coherent.
Semantic source replay can refute the lie only if trustworthy observed sources
distinguish these worlds; an ideal signature by itself does not supply them.
Trustworthy measurement, proof checking for faithfully encoded premises and
gateway mediation are possible assumptions/repairs, not proved properties of
arbitrary signers or organizations.
\end{proof}
